\section{Introduction}
For a nonempty finite alphabet $\calA$ with orthogonal command matrices $U_a$, a numerical experiment asks for one binary response after a word $w=a_1\cdots a_N$:
\begin{equation}\label{eq:target}
 \Pp(B=b\mid x,w,j)=\frac{1+b\rho\ip{e_j}{U_wx}}2,
 \qquad U_w=U_{a_N}\cdots U_{a_1},\quad b\in\{-1,1\}.
\end{equation}
Here $X\subset\Sphere^{D-1}$ is a nonempty finite seed set and $0<\rho\le1$ is fixed. A simulator receives the seed, reads every command once, and finally receives one coordinate query $j$. We count its largest persistent label alphabet, including a held binary answer. The simulator may use arbitrary stochastic rows, different rows at different epochs, and a different design for every prescribed horizon. Section~\ref{sec:model} gives the precise convention.

The distinction between a signed linear representation and a positive stochastic realization is classical. The mean in \eqref{eq:target} has signed dimension at most $D$. Positivity and common command updates impose different constraints. Equally importantly, a free external clock changes the relevant group. A single irrational rotation generates an infinite compact group, but its word of length $N$ is predetermined. Its entire effect can be absorbed in the decoder. The generated group alone cannot characterize bounded clocked width.

Let $R$ be the smallest invariant real subspace containing $X$, and restrict all commands to $R$. Put
\[
 H=\cl{\langle U_a|_R:a\in\calA\rangle}.
\]
The closure is a compact matrix group. Define its \emph{synchronization subgroup}
\begin{equation}\label{eq:sync-intro}
 \sync=\cl{\langle h(U_a^{-1}U_b)h^{-1}:a,b\in\calA,\ h\in H\rangle}.
\end{equation}
It is the least closed normal subgroup for which all commands have the same image in $H/\sync$. This group removes the common motion of words having the same length. It is defined on the active space, not on unused ambient coordinates.

\begin{theorem}[Boundedness dichotomy]\label{thm:main45}
For the experiment \eqref{eq:target}, the following are equivalent:
\begin{enumerate}
\item $\sync$ is finite;
\item $\sup_{N\ge1}W_{N,0}<\infty$.
\end{enumerate}
If $\sync$ is finite, deterministic command updates give
\[
 W_{N,0}\le\max\{2,|X|\,|\sync|\}\qquad(N\ge1).
\]
If $\sync$ is infinite, let $E=R\ominus R^{\sync^\circ}$ and
$\eta_X=\max_{x\in X}\norm{P_E x}$. Then $\eta_X>0$ and
\begin{equation}\label{eq:diverge45}
 W_{N,\varepsilon}\longrightarrow\infty
 \quad\text{for every fixed}\quad
 0\le\varepsilon<\frac{\rho\eta_X}{2\sqrt D}.
\end{equation}
No spectral gap, connectedness of $H$, or arithmetic hypothesis is required.
\end{theorem}

For symmetric coordinate seeds $X=\{\pm e_1,\ldots,\pm e_D\}$, the sufficient threshold in the infinite case is uniformly $\varepsilon<\rho/(2D)$. Theorem~\ref{thm:decision45} makes the exact boundedness dichotomy decidable when all input data lie in a specified real number field. It does not decide positivity of an infinite-dimensional action gap, nor determine the growth rate in every infinite case.

The proof identifies a finite-word obstruction before introducing probability. After removing a reference command, the sets of reachable vectors at equal lengths are nested. If their cardinalities remain bounded, their stabilized finite set is invariant under the synchronization subgroup. On the complement of its connected-component fixed space this is impossible. Compactness therefore gives a single finite packet whose every orbit has more than any proposed number of endpoint labels. The inherited endpoint-centroid contraction then forces a positive loss at each independent packet. This turns a finite-orbit criterion into a converse against arbitrary hidden stochastic states.

Two complementary realization statements are retained and made explicit from the concurrently published revision 44 \cite{GTF44}. Theorem~\ref{thm:flag45} gives the exact finite-horizon continuation-polytope formula for a prescribed width profile, including its actual terminal approximation error. Its ambient coordinates are all suffix responses; restricting them to the small signed representation would lose hidden realizations. Theorem~\ref{thm:comparison45} proves
\[
 \sup_{B\ge1}\frac{-\log(1-\Gamma_{\calA}(k,B))}{B}
                 \le \log\lambda_{\calA}(k,a),
\]
linking the previously separate packet and enclosure profiles. Equality is not asserted. The right side constructs vector-state realizations, whereas the continuation formula covers every hidden realization.

\paragraph{Provenance of the statements.}
The endpoint contraction, finite interval budget, circle optimization, and arithmetic results are retained from the predecessor development \cite{GTF35,GTF42,GTF43}, with full proofs below. The finite-horizon factorization theorem is a normalized controlled version of classical shift-cone realization, not a claim to have discovered positive realization \cite{BF,Ben22,Ohta}. The finite-bit compiler of revision 44 is also retained with its complete proof and code, and its scratch space remains a separate charged resource. The new organizing result is Theorem~\ref{thm:main45}, including removal of common clock motion, its arbitrary-hidden-state converse, and its finite algebraic decision consequence. The congruence bound used in that decision is classical and is proved explicitly. The packing--dilation balance of Corollary~\ref{cor:balance} remains a conditional corollary, not an exact quantitative classification of all alphabets. Section~\ref{sec:literature45} compares the optimization objects and approximation norms before the retained arithmetic applications.
